\documentclass[11pt]{report}
\usepackage{standalone}
\usepackage{amsmath,amssymb,amsfonts} % Typical maths resource packages
\usepackage{graphicx}                 % Packages to allow inclusion of graphics
\usepackage{color}                    % For creating coloured text and background
\usepackage[table]{xcolor}
\usepackage{hyperref}                 % For creating hyperlinks in cross references
\usepackage{algorithm}
%\usepackage{algorithmic}
\usepackage[noend]{algpseudocode}
\usepackage{longtable}
\usepackage{import}
\usepackage{lscape}
\usepackage{booktabs,colortbl}
\usepackage{tikz}
\usetikzlibrary{patterns}
\usetikzlibrary{trees}
\usepackage{forest}

\usepackage{pgfplots}
\usepackage{pgfplotstable}
\usepackage{endnotes}
\usepackage{cleveref}
\usepackage{caption}
\usepackage{subcaption}
\usepackage{afterpage}
\usepackage{comment}    
\usepackage{setspace}
\usepackage{listings}
\usepackage[margin=1in]{geometry}
\graphicspath{{./img/}}
\usepackage[toc,page]{appendix}
\usepackage{tocloft}
\usepackage{fancyhdr}

\pagestyle{fancy}

\fancyhf{}

\fancyheadoffset{0cm}

\renewcommand{\headrulewidth}{0pt} 

\renewcommand{\footrulewidth}{0pt}

\fancyhead[R]{\thepage}

\fancypagestyle{plain}{%

  \fancyhf{}%

  \fancyhead[R]{\thepage}%

}

\forestset{
  dir tree/.style={
    for tree={
      parent anchor=south west,
      child anchor=west,
      anchor=mid west,
      inner ysep=1pt,
      grow'=0,
      align=left,
      edge path={
        \noexpand\path [draw, \forestoption{edge}] (!u.parent anchor) ++(1em,0) |- (.child anchor)\forestoption{edge label};
      },
      font=\sffamily,
      if n children=0{}{
        delay={
          prepend={[,phantom, calign with current]}
        }
      },
      fit=band,
      before computing xy={
        l=2em
      }
    },
  }
}


% Define block styles

\tikzstyle{decision} = [diamond, draw, fill=blue!20, 
    text width=4.5em, text badly centered, node distance=3cm, inner sep=0pt]
\tikzstyle{block} = [rectangle, draw, fill=blue!20, 
    text width=5em, text centered, rounded corners, minimum height=4em]
\tikzstyle{line} = [draw, -latex']
\tikzstyle{cloud} = [draw, ellipse,fill=red!20, node distance=3cm,
    minimum height=2em]

\crefname{lstlisting}{listing}{listings}
\Crefname{lstlisting}{Listing}{Listings}
    

\newcommand\blankpage{%
    \null
    \thispagestyle{empty}%
    \addtocounter{page}{-1}%
    \newpage}

\newcommand{\bc}[1]{}

% A few result data files were not preserved in version control.  Where a plot
% cannot be redrawn, show an explicit placeholder rather than silently dropping
% the figure, so numbering, captions and cross-references stay intact.
\newcommand{\datalost}[1]{%
  \fbox{\parbox{0.42\linewidth}{\centering\footnotesize\vspace{3ex}
    Plot data not preserved\\ in this repository.\\[.5ex]
    \textit{#1}\vspace{3ex}}}%
}

% The Chapter 3 driver scripts (dfo/code) were never committed.  Fall back to a
% placeholder so the listing keeps its caption, label and cross-references and
% the document still builds.
\newcommand{\lstinputlistingif}[2][]{%
  \IfFileExists{#2}%
    {\lstinputlisting[#1]{#2}}%
    {\lstinputlisting[#1]{missing-listing.txt}}%
}


\newtheorem{theorem}{Theorem}[section]
\newtheorem{corollary}{Corollary}[theorem]
\newtheorem{lemma}[theorem]{Lemma}

\newenvironment{proof}[1][Proof]{\begin{trivlist}
\item[\hskip \labelsep {\bfseries #1}]}{\end{trivlist}}

\newcommand{\qed}{\nobreak \ifvmode \relax \else
      \ifdim\lastskip<1.5em \hskip-\lastskip
      \hskip1.5em plus0em minus0.5em \fi \nobreak
      \vrule height0.75em width0.5em depth0.25em\fi}

%\usetikzlibrary{external} 
%\tikzexternalize[prefix=tikz/]
%\tikzset{external/system call={lualatex --shell-escape -halt-on-error -interaction=batchmode -jobname ``\image'' ``\def\tikzexternalrealjob{main}\input{main}''}}
%\tikzset{external/export=false}



%\setlength{\textwidth}{15cm}
% The Graduate School required double spacing for the submitted copy.  For the
% reading copy (make reading) fall back to the tighter setting from header2.tex,
% which was the preferred layout before the formatting guidelines applied.
\ifdefined\readingcopy \linespread{1.1}\else \linespread{2}\fi
\parindent 1cm
\parskip 0.2cm
%\topmargin 1in %1cm  %.3cm
%\oddsidemargin 1in %1.5cm  %.7cm
%\evensidemargin 1in %1.5cm  %.7cm
%\textwidth 6.5in %13.5cm %15cm
%\textheight 9in %20cm %21cm


\newcommand{\bi}{\begin{itemize}}
\newcommand{\ei}{\end{itemize}}

\newcommand{\norm}[1]{ \left| \left| {#1} \right| \right|}
\newcommand{\magf}{ \left| f \right| }
\newcommand{\magp}{ \left| p \right| }
\newcommand{\magr}{ \left| r \right| }
\newcommand{\magomega}{ \left|  \omega  \right| }
\newcommand{\magE}{ \left| \mathcal{E} \right| }
\newcommand{\magG}{ \left| \mathcal{G} \right| }
\newcommand{\magL}{ \left| \mathcal{L} \right| }
\newcommand{\magM}{ \left| \mathcal{M} \right| }
\newcommand{\magT}{ \left| \mathcal{T} \right| }
\newcommand{\magV}{ \left| \mathcal{V} \right| }

\newcommand{\R}[1]{\mathbb{R}^{#1}}
\def\S{\mathbb{ S}}
\def\I{\mathbb{ I}}


\newcommand{\cA}{\mathcal{A}}
\newcommand{\cB}{\mathcal{B}}
\newcommand{\cD}{\mathcal{D}}
\newcommand{\cCALI}{\mathcal{CALI}}
\newcommand{\cDC}{\mathcal{DC}}
\newcommand{\cGR}{\mathcal{GR}}
\newcommand{\cE}{\mathcal{E}}
\newcommand{\cF}{\mathcal{F}}
\newcommand{\cG}{\mathcal{G}}
\newcommand{\cH}{\mathcal{H}}
\newcommand{\cI}{\mathcal{I}}
\newcommand{\cJ}{\mathcal{J}}
\newcommand{\cK}{\mathcal{K}}
\newcommand{\cL}{\mathcal{L}}
\newcommand{\cM}{\mathcal{M}}
\newcommand{\cN}{\mathcal{N}}
\newcommand{\cO}{\mathcal{O}}
\newcommand{\cS}{\mathcal{S}}
\newcommand{\cT}{\mathcal{T}}
\newcommand{\cU}{\mathcal{U}}
\newcommand{\cV}{\mathcal{V}}
\newcommand{\cW}{\mathcal{W}}
\newcommand{\cX}{\mathcal{X}}


\newcommand{\rd}{\pmb{d}}
\newcommand{\ry}{\pmb{y}}
\newcommand{\rxi}{\pmb{\xi}}
\newcommand{\romega}{\pmb{\omega}}
\newcommand{\rchi}{\pmb{\chi}}
\newcommand{\rdelm}{\pmb{\delta^m}}
\newcommand{\rye}{\pmb{y_e}}
\newcommand{\rdfe}{\pmb{\delta^y_e}}
\newcommand{\rdfee}{\pmb{\delta^y_{e_1}}}
\newcommand{\rdfeee}{\pmb{\delta^y_{e_2}}}
\newcommand{\rik}{\pmb{\delta^m_k}}
\newcommand{\rikk}{\pmb{\delta^m_{k_1}}}
\newcommand{\rikkk}{\pmb{\delta^m_{k_2}}}

\newcommand{\Expect}{\mathbb{E}}
\newcommand{\Exx}[1]{\mathbb{E}\left[ #1 \right]}
\newcommand{\Var}[1]{Var\left[ #1 \right]}
\newcommand{\CoVar}[2]{CoVar\left[ #1, #2 \right]}
\newcommand{\Prob}{\mathbb{P}}
\newcommand{\cvar}{\mathbb{CV}{\mathsf @}\mathbb{R}}
\newcommand{\grad}{\bigtriangledown}
\newcommand{\lb}{\left\{}
\newcommand{\rb}{\right\}}
\newcommand{\floor}[1]{\lfloor #1 \rfloor}

\newcommand{\defeq}{\stackrel{\rm def}{=}}



\makeatletter
\def\BState{\State\hskip-\ALG@thistlm}
\makeatother


\tikzset{
        hatch distance/.store in=\hatchdistance,
        hatch distance=10pt,
        hatch thickness/.store in=\hatchthickness,
        hatch thickness=2pt
    }

\makeatletter
    \pgfdeclarepatternformonly[\hatchdistance,\hatchthickness]{flexible hatch}
    {\pgfqpoint{0pt}{0pt}}
    {\pgfqpoint{\hatchdistance}{\hatchdistance}}
    {\pgfpoint{\hatchdistance-1pt}{\hatchdistance-1pt}}%
    {
        \pgfsetcolor{\tikz@pattern@color}
        \pgfsetlinewidth{\hatchthickness}
        \pgfpathmoveto{\pgfqpoint{0pt}{0pt}}
        \pgfpathlineto{\pgfqpoint{\hatchdistance}{\hatchdistance}}
        \pgfusepath{stroke}
    }


\usetikzlibrary{arrows,shapes,positioning,shadows,trees,svg.path}

\tikzset{
  basic/.style  = {draw, text width=2cm, drop shadow, font=\sffamily, rectangle},
  root/.style   = {basic, rounded corners=2pt, thin, align=center, fill=green!30, text width=2in},
  level 2/.style = {basic, rounded corners=6pt, thin,align=center, fill=green!45, text width=8em},
  oracle/.style = {basic, rounded corners=6pt, thin, fill=blue!30, text width=10em},
  level 2 oracle/.style = {basic, rounded corners=6pt, thin, fill=blue!10, text width=5.5em},
  level 3/.style = {basic, thin, align=left, fill=pink!60, text width=6.5em}
}


\makeindex
